% !TeX root = ../navier-stokes-manuscript.tex
\subsection{The temporal Tower}\label{sub:BodyFiniteTimeResponseEscalation}

Before \(T_*\), take the same solution on a strict interval where it remains classical and follow its Eulerian temporal responses at a fixed spatial coordinate:
\[
  u,
  \qquad
  \partial_tu,
  \qquad
  \partial_t^2u,
  \qquad
  \partial_t^3u,
  \qquad
  \ldots .
\]
Each row is the time derivative of the one before it.
They are successive responses of that field, not separate singularity scenarios and not material derivatives.
Write
\begin{equation}\label{eq:TemporalTowerDefinitionBody}
  V_n:=\partial_t^nu,
  \qquad
  Q_n:=\partial_t^np,
  \qquad
  n\in\Nzero.
\end{equation}
This notation fixes the family whose loss of control has just become consequential.

Navier--Stokes now tells us exactly what each row contains.
For \(V_0=u\),
\[
  V_1
  =\nu\Delta V_0
  -(V_0\cdot\nabla)V_0
  -\nabla Q_0.
\]
One time derivative can land on either copy of velocity in the transport term:
\[
  V_2
  =\nu\Delta V_1
  -(V_1\cdot\nabla)V_0
  -(V_0\cdot\nabla)V_1
  -\nabla Q_1.
\]
One more derivative produces the middle placement twice:
\[
  V_3
  =\nu\Delta V_2
  -(V_2\cdot\nabla)V_0
  -2(V_1\cdot\nabla)V_1
  -(V_0\cdot\nabla)V_2
  -\nabla Q_2.
\]
At order \(n\), the binomial coefficients count every way the time derivatives can divide between the transporting velocity and the velocity being transported:
\begin{equation}\label{eq:BodyTemporalVelocityRecurrence}
  V_{n+1}
  +\sum_{a=0}^{n}\binom{n}{a}(V_a\cdot\nabla)V_{n-a}
  +\nabla Q_n
  =\nu\Delta V_n,
  \qquad
  \nabla\cdot V_n=0.
\end{equation}
This recurrence is the exact Eulerian temporal Tower generated by Navier--Stokes.

Pressure appears in every row as well.
Taking the divergence of the original equation and using decay at spatial infinity recovers it from the velocity:
\begin{equation}\label{eq:PressureNormalization}
  -\Delta p=\partial_i\partial_j(u_i u_j),
  \qquad
  p=R_iR_j(u_i u_j).
\end{equation}
Here \(R_j:=\partial_j(-\Delta)^{-1/2}\), and the decaying normalization removes the arbitrary spatial constant in pressure at each time.
After \(n\) time derivatives,
\begin{equation}\label{eq:BodyTemporalPressureRecurrence}
  Q_n
  =R_iR_j
  \sum_{a=0}^{n}\binom{n}{a}
  (V_a)_i(V_{n-a})_j.
\end{equation}
The velocity rows determine \(Q_n\), and \(\nabla Q_n\) appears in the equation for \(V_{n+1}\).
The temporal Tower is already spatial inside the equation.
Viscosity contributes \(\Delta V_n\), transport contributes a spatial derivative, and pressure is recovered through spatial singular integrals of the same velocity products.
The resulting mixed coordinates are
\begin{equation}\label{eq:BodyMixedJetDefinition}
  J_{n,\alpha}:=\partial_t^n\partial_x^\alpha u,
  \qquad
  \Pi_{n,\alpha}:=\partial_t^n\partial_x^\alpha p,
  \qquad
  \alpha\in\Nzero^3.
\end{equation}
One spatial derivative splits the transport term again:
\[
  \partial_\ell\bigl((V_a\cdot\nabla)V_{n-a}\bigr)
  =((\partial_\ell V_a)\cdot\nabla)V_{n-a}
  +(V_a\cdot\nabla)\partial_\ell V_{n-a}.
\]
For a general multi-index \(\alpha\), Leibniz's rule gives
\begin{align}
  J_{n+1,\alpha}
  &+
  \sum_{a=0}^{n}
  \sum_{\beta\leq\alpha}
  \binom{n}{a}\binom{\alpha}{\beta}
  (J_{a,\beta}\cdot\nabla)
  J_{n-a,\alpha-\beta}
  +\nabla\Pi_{n,\alpha}
  \notag\\
  &=
  \nu\Delta J_{n,\alpha},
  \qquad
  \nabla\cdot J_{n,\alpha}=0.
  \label{eq:BodyMixedJetRecurrence}
\end{align}
The pressure coordinates satisfy
\begin{equation}\label{eq:BodyMixedPressureRecurrence}
  -\Delta\Pi_{n,\alpha}
  =
  \partial_x^\alpha\partial_i\partial_j
  \sum_{a=0}^{n}\binom{n}{a}
  (V_a)_i(V_{n-a})_j.
\end{equation}

Following the candidate singularity through successive responses led down the time direction, and the equation has opened spatial derivatives inside every rung.
Pressure stays tied to the same rows throughout.
The next count is the spatial height carried by each rung.
